% Generated draft: requires mathematical and source review.
\documentclass[11pt]{article}
\usepackage[margin=25mm]{geometry}
\usepackage{fontspec}
\setmainfont{Latin Modern Roman}
\usepackage{amsmath,amssymb,mathtools,unicode-math}
\setmathfont{Latin Modern Math}
\usepackage{enumitem,needspace,xurl,hyperref}
\setlength{\parindent}{0pt}
\setlength{\parskip}{5pt}
\setlength{\emergencystretch}{3em}
\newcommand{\N}{\mathbb{N}}
\newcommand{\Z}{\mathbb{Z}}
\newcommand{\Q}{\mathbb{Q}}
\newcommand{\R}{\mathbb{R}}
\newcommand{\C}{\mathbb{C}}
\newcommand{\F}{\mathbb{F}}
\newcommand{\A}{\mathbb{A}}
\newcommand{\PP}{\mathbb{P}}
\newcommand{\E}{\mathbb{E}}
\newcommand{\eps}{\varepsilon}
\newcommand{\dd}{\,\mathrm{d}}
\begin{document}
\section*{UnsolvedMath ID: 30004096}
\textbf{Source problem code:} catalogue entry
\subsection{Definitions and mathematical statement}
Let \(k\) be a field. The supplied record does not specify the permitted base fields; the formulation below is relative to \(k\), pending confirmation of the intended field hypotheses. A Hopf \(k\)-algebra is a unital associative \(k\)-algebra \(H\) equipped with a coassociative comultiplication \(\Delta:H\to H\otimes_k H\), a counit \(\varepsilon:H\to k\), and an antipode \(S:H\to H\), satisfying the bialgebra and antipode axioms. No bijectivity assumption on \(S\) is imposed here, since none is supplied.

Use the convention that \(H\) is noetherian if it satisfies the ascending chain condition on both left ideals and right ideals. The counit makes \(k\) a left and a right \(H\)-module by \(h\cdot a=\varepsilon(h)a\) and \(a\cdot h=a\varepsilon(h)\).

For this statement, use the ungraded, two-sided Artin--Schelter Gorenstein convention: \(H\) is \emph{AS-Gorenstein} if there is an integer \(d\geq 0\) such that the left and right regular modules have injective dimension
\[
\operatorname{injdim}_H({}_H H)=\operatorname{injdim}_{H^{\mathrm{op}}}(H_H)=d,
\]
and, for every integer \(i\geq 0\),
\[
\dim_k\operatorname{Ext}^{i}_H({}_H k,{}_H H)
=
\dim_k\operatorname{Ext}^{i}_{H^{\mathrm{op}}}(k_H,H_H)
=
\begin{cases}
1,&i=d,\\
0,&i\neq d.
\end{cases}
\]
Here right \(H\)-modules are viewed as left \(H^{\mathrm{op}}\)-modules. The one-dimensional conditions are conditions on \(k\)-vector spaces; they do not require the induced \(H\)-module structures on these Ext groups to be trivial.

\textbf{Question.} Is every noetherian Hopf \(k\)-algebra \(H\) AS-Gorenstein? Equivalently, under the conventions above, does every such \(H\) admit an integer \(d=d(H)\geq 0\) satisfying all the displayed conditions?

The source supplies neither base-field restrictions nor explicit sidedness and AS-Gorenstein conventions. The displayed formulation makes standard conventions explicit, but their agreement with the intended source formulation requires human mathematical review.
\subsection{Short English statement}
The question asks whether every noetherian Hopf algebra has finite injective dimension on both sides, with the Ext groups from the counit module vanishing except in one common degree, where each is one-dimensional. The base-field scope and precise conventions need source confirmation. The supplied record labels the problem open; that status is supplied dated curation, not an independently verified conclusion.
\subsection{Sources}
\begin{itemize}
\item Supplied local UnsolvedMath registry record, permanent ID \(30004096\), source problem code \texttt{catalogue entry}, title \emph{Artin--Schelter Gorenstein Property of Noetherian Hopf Algebras}, in the collection \emph{Oberwolfach Reports Open Problems}. No individual proposer is identified; the supplied proposed year is \(2019\).
\item Supplied source citation: \emph{Mini-Workshop: Cohomology of Hopf Algebras and Tensor Categories (2020)}, Oberwolfach Reports, pp.\ 663--693, DOI \texttt{10.4171/owr/2019/11}. \url{https://doi.org/10.4171/owr/2019/11}. The background describes Oberwolfach Reports as documenting mathematical workshops and their open-problem discussions.
\item Supplied literature source: \emph{Report on AS-Gorenstein Hopf algebras}. \url{https://arxiv.org/abs/2605.00774}.
\item Supplied literature source: \emph{Oberwolfach report}. \url{https://doi.org/10.4171/owr/2019/11}.
\item Literature assessment dated exactly \texttt{2026-08-22}: the supplied curation labels the universal assertion open, reports that many important classes are known, and states that no general proof or counterexample has been established. This is supplied dated curation derived from the OWR material, not independently verified here; its time-sensitive status and cited claims require independent verification.
\end{itemize}
\end{document}
